\uuid{QoU0}
\niveau{PCSI}
\module{Analyse}
\chapitre{Généralités sur les nombres réels}
%!TeX root=../../../encours.nouveau.tex
\duree{30}
\difficulte{2}
\auteur{Antoine Crouzet}
\datecreate{2024-12-01}
\titre{Caractérisation des intervalles}
\contenu{
\texte{Soit $I$ une partie non vide de $\R$ vérifiant la propriété suivante :
\[ \forall~x\in I,\,\forall~y\in I,\, x\leq y \implies \left[x,\,y\right] \subset I\]
L'objectif est de démontrer que $I$ est un intervalle.}
\question{On suppose que $I$ est majorée et non minorée.
	\begin{enumerate}
		\item Montrer qu'il existe un réel $a$ tel que $I~\subset \left]-\infty,\,a\right]$.
		\item Soit $x\in \left]-\infty,\,a\right[$. Montrer qu'il existe $y$ et $z$ dans $I$ tels que $y\leq x\leq z$.
		\item En déduire que $\left]-\infty,\,a\right[\subset I$ puis que $I$ est un intervalle.
	\end{enumerate}}
\reponse{
\begin{enumerate}
	\item $I$ est non vide (par hypothèse), et majorée. Elle admet une borne supérieure, que l'on note $a$. Ainsi, $I\subset \left]-\infty,\,a\right]$.
	\item Soit $x\in \left]-\infty,\,a\right[$. Puisque $I$ n'est pas minorée, il existe $y\in I$ tel que $y\leq x$ (sinon, $I$ est minorée par $y$).

	De plus, puisque $x<a$, par définition de la borne supérieure, il existe $z\in I$ tel que $x<z\leq a$.

	Ainsi, il existe bien deux éléments $y$ et $z$ de $I$ tels que $y\leq x\leq z$.
	\item Finalement, pour tout $x\in \left]-\infty,\,a\right[$, il existe $y$ et $z$ dans $I$ tels que $y\leq x\leq z$, c'est-à-dire $x\in \left[y,\,z\right]$. Or, par propriété de $I$, $\left[y,\,z\right]\subset I$, et donc $x\in I$.

	On a ainsi démontré que \[ \forall x\in \left]-\infty,\,a\right[,\quad x\in I \iff \left]-\infty,\,a\right[\subset I. \]
	Avec les deux inclusions démontrées en a) et c), on peut en déduire que $I=\left]-\infty,\,a\right[$ ou $I=\left]-\infty,\,a\right]$ : $I$ est bien un intervalle.
	\end{enumerate}
}
\question{Traiter de même les trois autres cas.}
\reponse{Les trois autres cas (majorée, minorée; minorée non majorée; non minorée non majorée) ce traite par le même principe.}
}
